\section{Related Work}
\label{sec:related}

\subsection{Parameter-Space Model Merging}
Model merging has emerged as a major paradigm for building unified, multi-task models without the steep computational overhead of joint training or multi-task fine-tuning \cite{wortsman2022model,ilharco2022editing}. Task Arithmetic \cite{ilharco2022editing} first demonstrated that subtracting a pre-trained base model's weights from a fine-tuned model's weights extracts a ``task vector'' that can be linearly scaled and added together to combine task capabilities. 
To mitigate cross-task parameter interference, several advanced weight merging heuristics have been developed. Ties-Merging \cite{yadav2023resolving} resolves sign conflicts and prunes low-magnitude parameter changes, while DARE \cite{yu2023language} uses randomized delta pruning to preserve the network's underlying scale. RegMean \cite{chohan2023regmean} optimizes merging based on closed-form covariance alignments of intermediate activations, and FisherMerge \cite{matena2022merging} utilizes diagonal Fisher Information matrices as weight-importance factors. Recent works like SyMerge \cite{kim2024symerge} introduce single-layer adapters to establish synergistic task manifolds, OrthoMerge \cite{spherelab2024orthomerge} utilizes orthogonal parameter mappings, and SAIM \cite{saim2026merge} focuses on isotropic sharpness-aware parameter-space exploration for continual learning. 

Despite their architectural and algorithmic elegance, all existing merging methods operate under the critical assumption of high-precision floating-point execution (FP16 or FP32). They are completely blind to low-bit quantization, failing to consider how parameter-space additions and sign-conflict resolutions warp weight distributions, making subsequent compression extremely challenging.

\subsection{Post-Training Quantization (PTQ)}
To deploy large models within tight VRAM and latency constraints, Post-Training Quantization (PTQ) has become a mandatory optimization step in production environments \cite{han2015deep,nagel2020up}. Standard PTQ aims to represent floating-point weights and activations using low-bit integers (typically INT4 or INT8) without retraining. Methods such as AdaRound \cite{nagel2020up} and BRECQ \cite{brecq2021nagel} optimize quantization rounding matrices using layer-reconstruction objectives. For transformer architectures, recent work has highlighted the severe impact of weight and activation outliers on quantization stability \cite{dettmers2022gpt3}. LLM.int8() \cite{dettmers2022gpt3} isolates extreme activation channels into a separate high-precision matrix multiplication stream. SmoothQuant \cite{xiao2023smoothquant} applies a mathematically equivalent scaling factor to migrate quantization difficulty from activations to weights, while AWQ \cite{lin2023awq} scales weights channel-wise based on activation magnitudes to protect salient channels. GPTQ \cite{frantar2022gptq} uses approximate second-order information to compensate for quantization errors in a block-wise fashion. To further compress models down to 2-4 bits, hybrid representations have been proposed: SqueezeLLM \cite{squeezellm2023} and SpQR \cite{spqr2023} decouple extreme weight outliers into high-precision sparse matrices while keeping the dense remainder in INT3/INT4 formats. 

However, existing PTQ frameworks are inherently single-task and single-model centric. They assume standard, smooth, single-domain weight distributions from standard training runs, and are completely blind to the heavy-tailed, scale-warped distributions that occur when multiple distinct task vectors are merged together in parameter space.

\subsection{The Merging-Quantization Gap}
As model merging and quantization have developed in separate, non-overlapping silos, deploying merged models on real-world hardware remains a major bottleneck. The only existing efforts to bridge this gap are limited to standard, uniform quantization of parameter soups, which suffer from severe precision loss at low bit-widths. In contrast, QP-Merge is the first framework that co-designs model merging and quantization. It explicitly addresses the joint heavy-tailed weight distribution of merged task vectors (via Outlier-Residual Decoupling) and corrects multi-task activation scale mismatches (via Quantization-Error Aware Scale Calibration) without requiring any expensive labeled target data. This pragmatic approach directly enables high-performance, low-cost edge deployment of multi-task models.
